%%%%%%%%%%%%%%%%%%%%%%%%%%%%%%%%%%%%%%%%%%%%%%%%%%%%%%%%%%%
% Lab report — Smart Cities RideSharing (Chicago)
% Template: rho-class v3 (Memo Jimenez, CC BY 4.0)
%%%%%%%%%%%%%%%%%%%%%%%%%%%%%%%%%%%%%%%%%%%%%%%%%%%%%%%%%%%

\documentclass[9pt,a4paper,twocolumn,twoside]{rho-class/rho}
\usepackage[english]{babel}

\doctype{Lab Report}
\title{Decoding the Intersection between Ridesourcing Dependency and Socio-Demographic Vulnerability in Chicago}

\author[{\textasteriskcentered},a,1]{Francesca Negri}
\author[{\textdagger},a,2]{Silvia Forcina Barrero}

\affil[a]{Data Science Lab on Smart Cities, University of Milan-Bicocca}

\dates{Compiled on \today}

\corres{\textsuperscript{\textasteriskcentered}Corresponding author: \href{mailto:francesca.negri@studenti.unimib.it}{francesca.negri@studenti.unimib.it}.}
\corres{\textsuperscript{\textdagger}Equal contribution.}

\docinfo{Analysis at Community Area level (77 CA, O'Hare excluded from socio-demographic modeling). Notebook: \texttt{SmartCities\_RideSharing.ipynb}.}

\journalname{DSLSC}
\journal{Smart Cities RideSharing}
\theday{\today}
\vol{1}
\no{1}

\begin{abstract}
    This report investigates spatial inequity in ridesourcing (Uber/Lyft) usage relative to socio-demographic vulnerability across Chicago's 77 Community Areas. We construct a Transport Usage Index (TUI) from Transportation Network Provider (TNP) trip data (2018--2024) and a composite Socio-Demographic Vulnerability Index (SDVI) from the Chicago Hardship Index and per capita income. Bivariate correlation analysis and side-by-side choropleth maps are used to test whether highly vulnerable areas depend more on ridesourcing.
\end{abstract}

\keywords{Ridesourcing, Smart Cities, Spatial Inequity, Chicago, Community Areas, Socio-Demographic Vulnerability}

\begin{document}

    \maketitle
    \thispagestyle{firststyle}

\section{Introduction}

    \rhostart{R}idesourcing platforms have reshaped urban mobility in large U.S. cities. In Chicago, uneven access to public transport and varying socio-demographic conditions raise a central policy question: does ridesourcing act as a mobility equalizer or as a costly substitute where transit is inadequate?

    We address this at the spatial scale of Chicago's 77 Community Areas (CA), comparing TNP usage intensity with composite vulnerability scores.

\section{Data Sources}

    \begin{itemize}
        \item \textbf{TNP trips} (Chicago Data Portal): monthly pickup counts by Community Area, 2018--2024.
        \item \textbf{Population} (CMAP Community Data Snapshots): resident population by CA and year.
        \item \textbf{Community Area boundaries} (Chicago Data Portal).
        \item \textbf{Hardship Index and per capita income} \cite{chicago_hardship}: socioeconomic indicators by neighborhood.
        \item \textbf{CCVI} \cite{chicago_ccvi}: auxiliary health vulnerability index (not included in SDVI).
    \end{itemize}

    Community Area 76 (O'Hare Airport) is excluded from socio-demographic modeling as a transport hub outlier.

\section{Methods}

    \subsection{Transport Usage Index (TUI)}

        TUI measures ridesourcing intensity per resident in Community Area $i$ during period $m$:
        \begin{equation}\label{eq:tui}
            TUI_{i,m} = \frac{\text{TNP pickups}_{i,m}}{\text{population}_{i,m}} \times 1000
        \end{equation}

    \subsection{Socio-Demographic Vulnerability Index (SDVI)}

        SDVI combines standardized Hardship Index and income disadvantage (lower per capita income implies higher vulnerability):
        \begin{equation}\label{eq:sdvi}
            SDVI_i = \frac{z(\text{Hardship}_i) + z(-\text{Income}_i)}{2}
        \end{equation}
        where $\text{Income}_i$ is per capita income in community area $i$. Higher SDVI values indicate greater socio-demographic vulnerability.

    \subsection{Correlation Analysis}

        We compute Pearson (linear) and Spearman (rank-based) correlations between annual mean TUI and SDVI across CAs. Significance is assessed at $\alpha = 0.05$; results are exploratory and do not imply causality.

\section{Results}

    \tabref{tab:indicators} summarizes the indicators computed in the notebook.

    \begin{table}[H]
        \caption{Indicators by Community Area}
        \label{tab:indicators}
        \begin{tabular}{
            >{\raggedright\arraybackslash}p{0.22\columnwidth}
            >{\raggedright\arraybackslash}p{0.62\columnwidth}
        }
        \toprule
        \textbf{Indicator} & \textbf{Description} \\
        \midrule
        TUI & TNP trips per 1{,}000 residents \\
        Hardship Index & Composite socioeconomic hardship score (1--100) \\
        Per capita income & Neighborhood-level income (USD) \\
        SDVI & Composite z-score mean of Hardship and income disadvantage \\
        CCVI & Auxiliary health/epidemiological vulnerability score \\
        \bottomrule
        \end{tabular}
        \tabletext{Insert correlation coefficients and map figures exported from the notebook.}
    \end{table}

    \begin{note}
        Export choropleth maps (TUI vs SDVI) and the scatter plot from the notebook and include them here as \texttt{figures/tui\_map.pdf} and \texttt{figures/sdvi\_map.pdf}.
    \end{note}

\section{Discussion}

    Interpret correlation signs in light of the research question. A positive association supports the \textit{substitute mobility} hypothesis; a negative one suggests concentration of ridesourcing in less vulnerable or central commercial areas. Spatial autocorrelation and land-use mix remain important confounders for future work.

\section{Code}

    Analysis is implemented in \texttt{SmartCities\_RideSharing.ipynb}. Helper functions live in \texttt{utils.py}.

    \begin{code}
        \lstinputlisting[language=python]{example.py}
        \label{lst:utils}
        \caption{Example Python utility (template). Replace with exported notebook snippets.}
    \end{code}

\section*{Acknowledgments}
\addcontentsline{toc}{section}{Acknowledgments}

    Course: Data Science Lab on Smart Cities, University of Milan-Bicocca.

\printbibliography

\end{document}
